\subsection{背景分析}

\subsubsection{外卖平台即时通信功能的普及与安全风险}

近年来，外卖平台已成为我国数字经济中规模最大、渗透率最高的生活服务场景之一。根据中国互联网络信息中心（CNNIC）发布的统计报告，截至2024年12月，我国网上外卖用户规模已达5.92亿人，占网民整体的53.4\%；至2025年12月，该规模进一步增长至6.30亿人，占网民整体的56.0\%\footnote{数据来源：中国互联网络信息中心（CNNIC）《中国互联网络发展状况统计报告》，中商产业研究院整理。}。中国饭店协会数据显示，2023年我国餐饮外卖市场规模约1.2万亿元，占餐饮收入的比重升至22.6\%\footnote{数据来源：中国饭店协会，转引自中国发展网2025年2月报道。}。外卖服务已从单纯的餐饮配送延伸至即时零售、药品配送、生鲜采购等多元场景，成为支撑城市生活运转的重要基础设施。

在庞大的用户基数和业务规模驱动下，外卖平台普遍引入了用户与骑手之间的即时通信（Instant Messaging）功能。该功能允许用户在订单创建后通过平台内置的聊天界面与配送骑手进行实时沟通，用于确认配送地址、沟通特殊需求、反馈配送问题等。然而，即时通信功能在提升配送效率的同时，也引入了新的安全隐患。与传统的电话沟通不同，订单即时通信产生的是可持久化存储的文本消息记录，这些记录中可能包含用户的家庭住址、生活习惯、消费偏好等敏感信息；同时，通信参与方的身份信息、订单信息与通信内容在系统中高度关联，一旦安全防护不足，将导致用户隐私的大面积暴露。

2025年4月，四川成都的王女士在外卖平台购买消毒液和计生用品后，配送骑手竟向其发送包含商品明细的骚扰短信\footnote{参见《法治日报》2025年5月8日第4版“法治经纬”栏目报道，新华社同步跟进报道。}。该事件表明，即便平台声称采用了虚拟号码加密技术，配送员仍可通过技术漏洞获取用户的真实手机号码及订单商品明细，并据此实施骚扰。类似事件并非个案：贵州省龙里县人民法院审结的一起案件中，某快递公司装车员在“每条信息0.8元”的利诱下连续偷拍快递面单，将数千条包含姓名、电话、住址及购买物品的隐私信息上传至网盘交易，最终4名涉案人员因侵犯公民个人信息罪获刑\footnote{参见贵州省龙里县人民法院相关案件审理报道，载于《法治日报》2025年5月8日。}。上述案例揭示了配送行业长期存在的隐私泄露隐患——在缺乏系统性安全设计的订单通信体系中，用户隐私处于事实上的“裸奔”状态。

\subsubsection{外卖订单通信中的四类安全问题}

通过对现有外卖平台订单通信流程的分析，并结合上述真实安全事件，我们归纳出外卖订单通信场景中存在的四类核心安全问题：

\textbf{（1）身份信息暴露}\quad 在传统外卖订单通信模式下，用户身份、订单身份与通信身份往往紧密绑定。骑手端在接单后可以直接获取用户的真实姓名、手机号码、收货地址等敏感信息，部分平台虽采用了虚拟号码技术，但据《法治日报》调查，实际运营中仍存在虚拟号短信通道未加密、多次通话后自动关联真实号码、配送端页面未完全屏蔽敏感字段等技术缺陷\footnote{参见《法治日报》2025年5月8日第4版调查报道。}。这些缺陷使得骑手能够在配送过程中获取超出业务必要范围的用户个人信息，进而可能引发骚扰、恐吓等二次侵害。

\textbf{（2）订单隐私泄露}\quad 现有外卖平台的即时通信功能主要关注业务连通性，缺乏对订单参与方合法性的系统化认证机制。在实际场景中，攻击者可能通过伪造订单身份、复用历史订单凭证或利用系统权限漏洞进入非自身订单的通信通道，从而窃取其他用户的通信内容或进行恶意骚扰。由于缺少基于密码学的订单认证机制，平台难以在不暴露用户真实身份的前提下验证通信参与方是否为合法的订单相关方。

\textbf{（3）聊天记录篡改}\quad 外卖订单通信产生的消息记录在发生配送纠纷、恶意投诉或法律诉讼时具有重要的证据价值。然而，现有平台的通信日志存储于平台中心化数据库中，既缺乏密码学完整性保护机制，也缺乏独立的第三方验证手段。当平台自身或具有数据库访问权限的内部人员篡改、删除或插入消息记录时，用户和监管机构难以发现和证明记录已被修改，从而导致通信记录的证据可信度不足。

\textbf{（4）纠纷追责困难}\quad 在完全匿名或弱匿名的通信环境中，用户和骑手之间的恶意行为（如骚扰、威胁、虚假投诉等）难以溯源到真实责任主体。如果系统默认暴露真实身份，则又会削弱隐私保护效果，形成“完全匿名不可管理”与“直接实名过度暴露”的两难困境。如何在保护日常通信隐私的同时保留异常情况下的条件溯源能力，是外卖订单通信安全需要解决的关键问题。

\subsubsection{国密算法推广与相关政策法规背景}

在外卖订单通信安全问题的技术层面，密码算法的选择至关重要。长期以来，RSA、AES、SHA-256 等国际通用密码算法在信息系统中被广泛使用，其安全性经过了较充分的公开分析和工程验证。然而，在我国商用密码推广和信息系统合规建设背景下，单纯依赖国际算法体系，可能在自主可控、国产密码适配、商用密码产品认证以及密码应用安全性评估等方面存在不足。为保障国家信息安全，我国大力推进商用密码算法的自主研发和推广应用。

2019年10月26日，第十三届全国人民代表大会常务委员会第十四次会议通过《中华人民共和国密码法》，自2020年1月1日起施行\footnote{《中华人民共和国密码法》，主席令第三十五号，2019年10月26日公布，自2020年1月1日起施行。}。该法明确将密码分为核心密码、普通密码和商用密码，规定商用密码用于保护不属于国家秘密的信息，鼓励商用密码技术的研究开发、学术交流、成果转化和推广应用。在密码法的框架下，国家密码管理局陆续发布了SM系列商用密码算法标准，其中SM3密码杂凑算法（GB/T 32905-2016）输出256 bit消息摘要，SM4分组密码算法（GB/T 32907-2016）采用128 bit分组长度和128 bit密钥长度，SM2椭圆曲线公钥密码算法（GB/T 32918）提供了数字签名、密钥交换和公钥加密等功能\footnote{GB/T 32905-2016《信息安全技术\ SM3密码杂凑算法》；GB/T 32907-2016《信息安全技术\ SM4分组密码算法》；GB/T 32918《信息安全技术\ SM2椭圆曲线公钥密码算法》。}。这些国密算法已在金融、电力、政务、交通等领域得到广泛应用，为我国信息系统的安全自主可控提供了技术基础。

在个人信息保护层面，2021年8月20日通过的《中华人民共和国个人信息保护法》（自2021年11月1日起施行）确立了处理个人信息应遵循的“合法、正当、必要和诚信原则”，并明确提出了“最小必要原则”——处理个人信息应当具有明确、合理的目的，并应当与处理目的直接相关，采取对个人权益影响最小的方式\footnote{《中华人民共和国个人信息保护法》第五条、第六条。}。对于外卖配送场景而言，配送员仅需知晓收件人地址和虚拟号码，商品明细、真实电话号码等均属非必要信息，不应向配送环节披露\footnote{参见《法治日报》2025年5月8日第4版，中国传媒大学程科副教授及北京嘉潍律师事务所赵占领律师的专家意见。}。此外，《中华人民共和国数据安全法》（自2021年9月1日起施行）要求数据处理者建立健全全流程数据安全管理制度，采取相应的技术措施保障数据安全\footnote{《中华人民共和国数据安全法》第二十七条。}。国务院颁布的《“十四五”数字经济发展规划》亦提出要加强数据加密技术发展，提升数字经济安全体系的保障能力\footnote{《“十四五”数字经济发展规划》，国务院2021年12月印发。}。

上述法律法规和政策文件为外卖订单通信安全系统的设计提供了明确的合规指引：系统应当在“最小必要原则”的指导下，采用国密算法构建身份匿名、订单认证、消息保护和日志审计机制，在保护用户隐私的同时满足平台管理和纠纷追责的监管要求。
